\documentclass[12pt]{article}
\usepackage[utf8]{inputenc}
\usepackage[T1]{fontenc}
\usepackage[margin=1in]{geometry}
\usepackage{setspace}
\usepackage{xcolor}
\usepackage[colorlinks=true, urlcolor=blue!60!black, linkcolor=blue!60!black]{hyperref}
\setlength{\parindent}{2em}
\setlength{\parskip}{1.25em}
\renewcommand{\baselinestretch}{1.0}
\begin{document}
\begin{center}
{\LARGE \textit{Jacobson v. Massachusetts}, 197 U.S. 11 (1905)\par}
\vspace{0.8em}
{\large Case Brief\par}
\end{center}
\vspace{1em}

\section{Memory Jogger}

In the midst of a 1902 Cambridge smallpox outbreak, the Supreme Court upheld Massachusetts's compulsory adult vaccination statute and its \$5 fine as a valid exercise of the state police power, holding that Fourteenth Amendment liberty is ``liberty regulated by law'' subject to reasonable restraints for the common good (\href{https://www.courtlistener.com/opinion/96230/jacobson-v-massachusetts/}{CourtListener}; \href{https://supreme.justia.com/cases/federal/us/197/11/}{Justia}).

\section{Detailed Case Facts}

Revised Laws of Massachusetts, chapter 75, \S~137, authorized the board of health of any city or town, if ``in its opinion, it is necessary for the public health or safety,'' to require and enforce the vaccination and revaccination of all inhabitants while providing them with free vaccination (\href{https://supreme.justia.com/cases/federal/us/197/11/}{Justia}). Any person over twenty-one years of age and not under guardianship who refused or neglected to comply forfeited \$5, while \S~139 created an explicit medical exemption only for ``children who present a certificate, signed by a registered physician that they are unfit subjects for vaccination'' (\href{https://supreme.justia.com/cases/federal/us/197/11/}{Justia}).

On February 27, 1902, citing the prevalence and continuing spread of smallpox in Cambridge, the Cambridge Board of Health ordered the vaccination or revaccination of all city inhabitants who had not been successfully vaccinated since March 1, 1897, and subsequently empowered a designated physician to enforce the order (\href{https://supreme.justia.com/cases/federal/us/197/11/}{Justia}). Henning Jacobson, a Swedish Lutheran pastor residing in Cambridge who was over twenty-one and not under guardianship, refused the Board chairman's offer of free vaccination (\href{https://supreme.justia.com/cases/federal/us/197/11/}{Justia}; \href{https://en.wikipedia.org/wiki/Jacobson_v._Massachusetts}{Wikipedia}). Jacobson maintained that he had suffered severe, prolonged illness following a childhood vaccination in Sweden and had witnessed a similar adverse reaction in his son, leading him to fear that vaccination was dangerous and ineffective (\href{https://supreme.justia.com/cases/federal/us/197/11/}{Justia}; \href{https://pmc.ncbi.nlm.nih.gov/articles/PMC1449224/}{Mariner et al., Am.\ J.\ Pub.\ Health}). Crucially, however, Jacobson did not offer medical proof that, in his then-present adult condition in 1902, he was medically unfit for vaccination (\href{https://supreme.justia.com/cases/federal/us/197/11/}{Justia}).

\section{Procedural History}

On July 17, 1902, Massachusetts brought a criminal complaint against Jacobson in a local Massachusetts inferior court for refusing to comply with the Cambridge Board of Health regulation (\href{https://supreme.justia.com/cases/federal/us/197/11/}{Justia}). After Jacobson appealed an initial conviction and \$5 fine to the Superior Court of Middlesex County, he pleaded not guilty and sought to introduce fourteen offers of proof contesting the safety and efficacy of vaccination and recounting his and his son's past adverse reactions (\href{https://supreme.justia.com/cases/federal/us/197/11/}{Justia}; \href{https://purposedrivenlawyers.com/wp-content/uploads/2024/06/Jacobson-v-Commonwealth-of-Massachusetts.pdf}{U.S.\ Reports / Westlaw Reprint}). The trial judge excluded all fourteen offers of proof as immaterial and rejected Jacobson's requested jury instructions asserting that \S~137 violated the Preamble, ``spirit,'' and Fourteenth Amendment (Privileges or Immunities, Due Process, and Equal Protection Clauses) of the U.S.\ Constitution (\href{https://supreme.justia.com/cases/federal/us/197/11/}{Justia}). The jury returned a guilty verdict.

The trial court continued the case for the opinion of the Supreme Judicial Court of Massachusetts, which overruled all of Jacobson's exceptions in \textit{Commonwealth v.\ Pear} / \textit{Commonwealth v.\ Jacobson}, 183 Mass.\ 242, 66 N.E.\ 719 (1903) (\href{https://supreme.justia.com/cases/federal/us/197/11/}{Justia}; \href{https://caselaw.findlaw.com/court/us-supreme-court/197/11.html}{FindLaw}). Pursuant to the verdict, the Superior Court sentenced Jacobson to pay a \$5 fine and ordered him committed until the fine was paid (\href{https://supreme.justia.com/cases/federal/us/197/11/}{Justia}). Jacobson brought a writ of error to the U.S.\ Supreme Court (No.\ 70), which heard argument on December 6, 1904, and issued its decision on February 20, 1905 (\href{https://www.oyez.org/cases/1900-1940/197us11}{Oyez}; \href{https://tile.loc.gov/storage-services/service/ll/usrep/usrep197/usrep197011/usrep197011.pdf}{Library of Congress}).

\section{Judicial Votes}

\textbf{Vote:} 7--2, affirming the judgment of the Superior Court of Massachusetts (\href{https://supreme.justia.com/cases/federal/us/197/11/}{Justia}; \href{https://www.oyez.org/cases/1900-1940/197us11}{Oyez}).

\begin{itemize}
\item \textbf{Majority (7):} Justice John Marshall Harlan (author), joined by Chief Justice Melville Weston Fuller and Associate Justices Henry Billings Brown, Edward Douglass White, Joseph McKenna, Oliver Wendell Holmes, Jr., and William Rufus Day.
\item \textbf{Concurrences:} None.
\item \textbf{Dissents (2):} Associate Justices David Josiah Brewer and Rufus Wheeler Peckham dissented without written opinion (\href{https://supreme.justia.com/cases/federal/us/197/11/}{Justia}).
\end{itemize}

\section{Holding}

\textbf{Judgment:} Affirmed. Massachusetts's compulsory vaccination statute, as applied to a healthy adult during a smallpox epidemic, did not violate the Fourteenth Amendment (\href{https://supreme.justia.com/cases/federal/us/197/11/}{Justia}).

\textbf{Rule Announced:} Under its reserved police power, a state may enact reasonable public health regulations---including delegating authority to local boards of health to require adult vaccination during a disease outbreak---so long as the measure does not contravene the Federal Constitution. Judicial invalidation of a purported public-health statute is warranted only if the law ``has no real or substantial relation to those objects, or is, beyond all question, a plain, palpable invasion of rights secured by the fundamental law'' (\href{https://supreme.justia.com/cases/federal/us/197/11/}{Justia}). Moreover, exempting medically unfit children without an express statutory exemption for adults does not violate Equal Protection where the law applies equally to all adults in like condition (\href{https://supreme.justia.com/cases/federal/us/197/11/}{Justia}).

\textbf{Holding vs.\ Dicta:} The \emph{holding} is deliberately narrow: the Court decided ``only that the statute covers the present case''---an adult who appeared in perfect health and a fit subject for vaccination refusing compliance during an active smallpox outbreak in Cambridge---and that the Preamble and general ``spirit'' of the Constitution confer no independent substantive rights (\href{https://supreme.justia.com/cases/federal/us/197/11/}{Justia}). Key \emph{dicta} include Justice Harlan's extended discussion of social-compact theory, military conscription, and maritime quarantine, as well as his closing cautionary reservation that courts would be competent to intervene if a state attempted to enforce vaccination against an adult for whom it could be shown with reasonable certainty that vaccination would ``seriously impair his health or probably cause his death'' (\href{https://supreme.justia.com/cases/federal/us/197/11/}{Justia}).

\section{Analysis of Opinions}

\subsection{Majority Opinion (Justice Harlan)}

Justice Harlan began by clearing away Jacobson's threshold textual arguments. The Constitution's Preamble indicates general purposes but confers no substantive powers or enforceable individual rights, and while the ``spirit'' of the Constitution is respected, it ``is to be collected chiefly from its words'' (\href{https://supreme.justia.com/cases/federal/us/197/11/}{Justia}). Accepting the Massachusetts Supreme Judicial Court's construction of \S~137 as binding, Harlan framed the core substantive due process issue: whether compulsory vaccination under local board authority exceeded Massachusetts's police power or invaded Fourteenth Amendment liberty.

Harlan grounded state authority in the police power---retained upon entering the Union and encompassing ``quarantine laws and `health laws of every description'''---and in the social compact recognized in the Massachusetts Constitution of 1780 (\href{https://supreme.justia.com/cases/federal/us/197/11/}{Justia}). Rejecting an absolutist view of bodily autonomy, Harlan reasoned that constitutional liberty does not entitle any individual to be ``at all times and in all circumstances, wholly freed from restraint''; rather, ``[t]here are manifold restraints to which every person is necessarily subject for the common good,'' including quarantine of healthy travelers and military conscription (\href{https://supreme.justia.com/cases/federal/us/197/11/}{Justia}). On the principle of ``self-defense, of paramount necessity,'' a community may protect itself against an epidemic (\href{https://supreme.justia.com/cases/federal/us/197/11/}{Justia}).

Turning to the standard of review, Harlan rejected Jacobson's request to have a jury second-guess medical science. Because medical consensus and international experience strongly supported vaccination---a ``common belief'' of which courts take judicial notice---choosing among competing medical theories belonged to the legislature, not a court or jury (\href{https://supreme.justia.com/cases/federal/us/197/11/}{Justia}). Citing \textit{Mugler v.\ Kansas}, 123 U.S.\ 623 (1887), Harlan articulated a deferential two-prong test: courts may strike down a police-power health regulation only when it bears ``no real or substantial relation'' to public health or safety, or is ``beyond all question, a plain, palpable invasion of rights secured by the fundamental law'' (\href{https://supreme.justia.com/cases/federal/us/197/11/}{Justia}). He also rejected the Equal Protection challenge because ``obvious reasons'' distinguished children of tender years from adults, and the law treated all adults equally (\href{https://supreme.justia.com/cases/federal/us/197/11/}{Justia}). Finally, Harlan underscored judicial guardrails: police power cannot be exercised in an ``arbitrary and oppressive'' manner, and applying the statute to an adult whose medical condition rendered vaccination life-threatening would be ``cruel and inhuman in the last degree,'' warranting judicial intervention (\href{https://supreme.justia.com/cases/federal/us/197/11/}{Justia}).

\subsection{Dissenting Opinions (Justices Brewer and Peckham)}

Justices David Josiah Brewer and Rufus Wheeler Peckham dissented without publishing written opinions (\href{https://supreme.justia.com/cases/federal/us/197/11/}{Justia}). Their silent dissents are nonetheless doctrinally revealing in the context of \textit{Lochner}-era substantive due process: just two months later, Justice Peckham authored \textit{Lochner v.\ New York}, 198 U.S.\ 45 (1905)---joined by Justice Brewer---striking down a New York bakery-hours statute as an unconstitutional interference with liberty of contract, whereas Justice Harlan dissented in \textit{Lochner} by invoking \textit{Jacobson}'s deferential police-power framework (\href{https://pmc.ncbi.nlm.nih.gov/articles/PMC1449223/}{Gostin, Am.\ J.\ Pub.\ Health}; \href{https://www.cambridge.org/core/books/covid19-and-the-law/jacobson-question/A463868C48886F30BD127463D9796626}{Cambridge Univ.\ Press}).

\section{Examples: Future Applications}

\begin{enumerate}
\item \textbf{Same-Side (Upheld): State School-Entry Immunization Mandate.} A state statute bars unvaccinated children from attending public or private schools unless they qualify for a medical exemption. This falls squarely on the same side as \textit{Jacobson} because the state is exercising its police power to curb communicable disease through a reasonable, non-arbitrary condition bearing a real and substantial relation to public health (indeed, the Supreme Court applied \textit{Jacobson} to uphold such an ordinance in \textit{Zucht v.\ King}, 260 U.S.\ 174 (1922)) (\href{https://supreme.justia.com/cases/federal/us/197/11/}{Justia}; \href{https://historyofvaccines.org/blog/faith-law-and-needle-history-religious-exemptions-vaccination-united-states/}{History of Vaccines}).

\item \textbf{Same-Side (Upheld): Short-Term Quarantine of Exposed Travelers.} During an outbreak of a highly contagious hemorrhagic fever, a state health department isolates passengers arriving from an affected region for fourteen days at a monitored facility until medical screening confirms they are non-infectious. This tracks \textit{Jacobson}'s explicit approval of maritime quarantine for yellow fever or cholera: the restraint is temporary, tailored to an imminent epidemic threat, and grounded in community self-defense (\href{https://supreme.justia.com/cases/federal/us/197/11/}{Justia}).

\item \textbf{Opposite-Side (Invalidated): Enforcing a Vaccine Mandate Against a Medically Contraindicated Adult.} A city board of health fines and threatens to jail an adult with documented severe anaphylaxis to vaccine components, refusing to recognize any adult medical exemption despite unanimous physician certification that the shot will likely kill him. This falls on the opposite side under \textit{Jacobson}'s own closing reservation: compelling vaccination when it can be shown with reasonable certainty to ``seriously impair his health or probably cause his death'' is ``cruel and inhuman in the last degree'' and arbitrary and oppressive as applied (\href{https://supreme.justia.com/cases/federal/us/197/11/}{Justia}).

\item \textbf{Opposite-Side (Invalidated): Discriminatory Pandemic Restrictions Singling Out Houses of Worship.} During a respiratory virus surge, a governor caps attendance at religious services at ten persons regardless of sanctuary size while allowing comparable secular venues (such as large retail stores and acupuncture facilities) to operate without numerical caps. Under modern constitutional doctrine reaffirmed in \textit{Roman Catholic Diocese of Brooklyn v.\ Cuomo}, 592 U.S.\ 14, 141 S.\ Ct.\ 63 (2020), \textit{Jacobson} does not suspend heightened scrutiny for enumerated First Amendment rights or authorize regulations that single out religion for harsher treatment (\href{https://www.plainsite.org/opinions/1jlepwj8p/roman-catholic-diocese-of-brooklyn-v-cuomo/}{PlainSite}; \href{https://harvardlawreview.org/print/vol-134/constitutional-constraints-on-free-exercise-analogies/}{Harvard Law Review}).

\item \textbf{Fence-Sitter (Close Case): Indefinite Statewide Adult Vaccine Mandate for a Low-Fatality Endemic Virus.} A state legislature enacts a permanent \$500-per-month fine for any adult who declines an annual booster for an endemic respiratory virus with a low mortality rate and vaccines that reduce severe illness more than transmission. Supporters would argue under \textit{Jacobson} that legislatures, not courts, weigh competing epidemiological evidence under the ``real or substantial relation'' test (\href{https://supreme.justia.com/cases/federal/us/197/11/}{Justia}). Challengers would argue that \textit{Jacobson} turned on an acute, high-mortality epidemic (``paramount necessity''), a modest \$5 fine, and a vaccine preventing community transmission---whereas an indefinite, escalating financial penalty for a less transmissible-blocking intervention tests \textit{Jacobson}'s proportionality and necessity limits as well as modern bodily-integrity doctrine under \textit{Cruzan v.\ Director, Missouri Department of Health}, 497 U.S.\ 261 (1990) (\href{https://www.law.cornell.edu/constitution-conan/amendment-14/right-to-refuse-medical-treatment-and-substantive-due-process}{Cornell LII}; \href{https://pmc.ncbi.nlm.nih.gov/articles/PMC1449223/}{Gostin, Am.\ J.\ Pub.\ Health}).
\end{enumerate}

\section{Critique}

\textbf{Originalist, Textualist, and Libertarian Perspectives:} Originalist and libertarian scholars---as well as Justice Neil Gorsuch concurring in \textit{Roman Catholic Diocese of Brooklyn v.\ Cuomo}, 141 S.\ Ct.\ 63, 70 (2020)---argue that \textit{Jacobson} was a modest, pre-incorporation substantive due process decision applying early rational-basis review to an unenumerated bodily-autonomy claim backed only by a \$5 fine, not a ``super-precedent'' suspending the Bill of Rights during emergencies (\href{https://digitalcommons.law.buffalo.edu/cgi/viewcontent.cgi?article=4934&context=buffalolawreview}{Blackman, Buffalo L.\ Rev.}; \href{https://www.cambridge.org/core/books/covid19-and-the-law/jacobson-question/A463868C48886F30BD127463D9796626}{Cambridge Univ.\ Press}). Josh Blackman critiques the ``irrepressible myth'' of \textit{Jacobson}, warning that courts anachronistically grafted a 1905 \textit{Lochner}-era police-power opinion onto modern tiered-scrutiny doctrine to create an executive ``escape hatch'' during COVID-19 (\href{https://digitalcommons.law.buffalo.edu/cgi/viewcontent.cgi?article=4934&context=buffalolawreview}{Blackman, Buffalo L.\ Rev.}). Critics across the spectrum also emphasize \textit{Jacobson}'s darkest doctrinal legacy: Justice Holmes's citation of \textit{Jacobson} as the sole precedent supporting compulsory eugenic sterilization in \textit{Buck v.\ Bell}, 274 U.S.\ 200, 207 (1927), illustrating the danger of unchecked deference to purported scientific consensus (\href{https://pmc.ncbi.nlm.nih.gov/articles/PMC1449224/}{Mariner et al., Am.\ J.\ Pub.\ Health}; \href{https://www.independent.org/tir/2021-spring/infectious-diseases-and-government-growth/}{Independent Institute}).

\textbf{Progressive and Public-Health Doctrinal Perspectives:} Public-health scholars such as Lawrence Gostin view \textit{Jacobson} more favorably as a foundational counterweight to \textit{Lochner}-era laissez-faire activism, establishing a balanced constitutional floor built on four principles: public-health necessity, reasonable means, proportionality, and harm avoidance (\href{https://pmc.ncbi.nlm.nih.gov/articles/PMC1449223/}{Gostin, Am.\ J.\ Pub.\ Health}). Yet progressive civil-liberties scholars such as Wendy Mariner, George Annas, and Leonard Glantz caution that \textit{Jacobson} predates modern substantive due process, Bill of Rights incorporation, and strict scrutiny for fundamental rights; consequently, it should be read as affirming democratic authority to address genuine public-health threats rather than as licensing arbitrary intrusions on bodily integrity or minority rights (\href{https://pmc.ncbi.nlm.nih.gov/articles/PMC1449224/}{Mariner et al., Am.\ J.\ Pub.\ Health}).

\textbf{Underlying Competing Values:} \textit{Jacobson} exposes the enduring constitutional tension between (1) individual bodily autonomy and freedom from majoritarian coercion, (2) collective survival and the social compact's ``common good,'' and (3) institutional competence---specifically, how much courts should defer to democratically accountable legislatures and public-health experts during scientific uncertainty without abdicating their duty to enforce constitutional boundaries (\href{https://www.cambridge.org/core/books/covid19-and-the-law/jacobson-question/A463868C48886F30BD127463D9796626}{Cambridge Univ.\ Press}).

\section{Key Quotations}

\begin{quote}
``But the liberty secured by the Constitution of the United States to every person within its jurisdiction does not import an absolute right in each person to be, at all times and in all circumstances, wholly freed from restraint. There are manifold restraints to which every person is necessarily subject for the common good.''
\end{quote}
\noindent --- Justice Harlan, \textit{Jacobson v.\ Massachusetts}, 197 U.S.\ 11, 26 (1905) (\href{https://supreme.justia.com/cases/federal/us/197/11/}{Justia}).

\begin{quote}
``Upon the principle of self-defense, of paramount necessity, a community has the right to protect itself against an epidemic of disease which threatens the safety of its members.''
\end{quote}
\noindent --- Justice Harlan, \textit{Jacobson v.\ Massachusetts}, 197 U.S.\ 11, 27 (1905) (\href{https://supreme.justia.com/cases/federal/us/197/11/}{Justia}).

\begin{quote}
``[I]f a statute purporting to have been enacted to protect the public health, the public morals, or the public safety has no real or substantial relation to those objects, or is, beyond all question, a plain, palpable invasion of rights secured by the fundamental law, it is the duty of the courts to so adjudge, and thereby give effect to the Constitution.''
\end{quote}
\noindent --- Justice Harlan, \textit{Jacobson v.\ Massachusetts}, 197 U.S.\ 11, 31 (1905) (quoting \textit{Mugler v.\ Kansas}, 123 U.S.\ 623, 661 (1887)) (\href{https://supreme.justia.com/cases/federal/us/197/11/}{Justia}).

\begin{quote}
``[T]he police power of a State, whether exercised by the legislature or by a local body acting under its authority, may be exerted in such circumstances or by regulations so arbitrary and oppressive in particular cases as to justify the interference of the courts to prevent wrong and oppression.''
\end{quote}
\noindent --- Justice Harlan, \textit{Jacobson v.\ Massachusetts}, 197 U.S.\ 11, 38 (1905) (\href{https://supreme.justia.com/cases/federal/us/197/11/}{Justia}).


\end{document}
